%! Bias in Samples and Surveys — Unit 1, Lesson 1.4 — Slides (Beamer)
% PILOT SHAPE (2026-09-06): modeled on AP Statistics lesson 1.4 minus its AP Practice page.
% GRADUAL RELEASE deck: title → targets → warm-up → hook (unresolved) → I DO divider → two or
% three frames per notes section → WE DO (a live surface, un-answered) → spoken debrief → close &
% START THE HOMEWORK, which is the individual practice. No group activity, no practice launch,
% no exit ticket. There is no spoiler rule: the notes frames ARE the direct instruction.
\documentclass[aspectratio=169, 11pt]{beamer}
\usepackage{saar-beamer}   % provides \ceruleanheader{} and \sectionlabel[color]{}

\newcommand{\redink}[1]{{\color{redacc}\bfseries #1}}

\begin{document}

% TITLE SLIDE (CourseName is not defined in beamer, so the name is literal)
{
\setbeamercolor{background canvas}{bg=cerulean}
\begin{frame}[plain]
  \vspace{2.8cm}\hspace{0.55cm}%
  \begin{minipage}{0.9\textwidth}
    {\color{goldacc}\bfseries\small LESSON 1.4}\\[0.5cm]
    {\color{white}\bfseries\LARGE Bias in Samples and Surveys}\\[1.2cm]
    {\color{frostmid}\small Statistical Analysis \& Algebraic Reasoning~$\cdot$~Unit 1}
  \end{minipage}
\end{frame}
}

% ── LEARNING TARGETS ─────────────────────────────────────────────────────────
\begin{frame}[t, plain]
  \ceruleanheader{Today's targets}
  \sectionlabel{Lesson 1.4}
  By the end of today, I can\ldots
  \begin{itemize}\setlength{\itemsep}{5pt}
    \item explain what makes a study \textbf{biased} and tell bias apart from \textbf{sampling
          variability} --- a direction, not bad luck.
    \item identify \textbf{sampling bias} in a plan --- \textbf{undercoverage} and
          \textbf{nonresponse} --- by asking \emph{could that person have been chosen?}
    \item name the forms of \textbf{response bias}, including \textbf{question order bias}, and
          say which \textbf{direction} a plan pushes the estimate.
    \item explain why a \textbf{larger sample does not fix bias}, and describe a change that
          removes it.
  \end{itemize}
  \begin{block}{How today works}
    Warm-up $\to$ \textbf{notes together} $\to$ \textbf{one survey we diagnose as a class} $\to$
    debrief $\to$ \textbf{you start the homework here, alone}.
  \end{block}
\end{frame}

% ── WARM-UP ──────────────────────────────────────────────────────────────────
\begin{frame}[t, plain]
  \ceruleanheader{Warm-Up --- 5 minutes, silent}
  \sectionlabel{Harbor Point Community Pool: 1,500 members}
  \vspace{0.3cm}
  \begin{center}
    \begin{minipage}{0.86\textwidth}
      \begin{enumerate}\setlength{\itemsep}{7pt}
        \item The office's records: $300$ of the $1{,}500$ members mainly swim laps. What percent?
        \item The study we threw out last class: $27$ of the $75$ asked on a Tuesday evening said
              laps. What percent of those $75$?
        \item Use the $36\%$ to predict the whole membership. How many members? The records say
              $300$ --- off by how many, and too high or too low?
      \end{enumerate}
    \end{minipage}
  \end{center}
  \vspace{0.4cm}
  \begin{center}
    {\color{cerulean}\bfseries Algebra 1 arithmetic. Keep the page --- section 1 opens on these
    four numbers.}
  \end{center}
\end{frame}

% ── HOOK — leave it UNRESOLVED; the second half of section 2 settles it ──────
{
\setbeamercolor{background canvas}{bg=cerulean}
\begin{frame}[plain]
  \vspace{1.2cm}
  \begin{center}
    {\color{frostmid}\bfseries\small THE TUESDAY-EVENING STUDY SAID 36\%. THE RECORDS SAY 20\%.}\\[0.7cm]
    {\color{white}\bfseries\Large The staff member says the problem was size. She will run the
    same Tuesday-evening survey for ten weeks and ask 750 members instead of 75.}\\[0.9cm]
    {\color{goldacc}\bfseries\large Will her new number land closer to the truth?}\\[0.5cm]
    {\color{frostmid}\small Vote: yes, no, or it depends. Write one line of reason in the hook
    box. We settle this at the end of section 2.}
  \end{center}
\end{frame}
}

% ── I DO — direct instruction begins ─────────────────────────────────────────
{
\setbeamercolor{background canvas}{bg=cerulean}
\begin{frame}[plain]
  \vspace{1.8cm}\hspace{0.8cm}%
  \begin{minipage}{0.86\textwidth}
    {\color{goldacc}\bfseries\small GUIDED NOTES \ \textbullet\ \ I DO}\\[0.7cm]
    {\color{white}\bfseries\Large Open to your Guided Notes page. Fill each blank as we
    name it --- not before.}\\[0.9cm]
    {\color{frostmid}\small Five terms go in the vocabulary box today. Write each one the
    moment we use it.}
  \end{minipage}
\end{frame}
}

% ── 1a. BIAS IS A DIRECTION ──────────────────────────────────────────────────
\begin{frame}[t, plain]
  \ceruleanheader{Bias is a direction, not bad luck}
  \sectionlabel{notes section 1 $\cdot$ variability scatters around the truth; bias leans}
  \vspace{2pt}
  The office's records: $20\%$ of all $1{,}500$ members mainly swim laps. Three samples of $75$, two ways:
  \begin{center}
    \small
    \renewcommand{\arraystretch}{1.3}
    \begin{tabular}{@{} l c c c c @{}}
    \hline
    \textbf{How the 75 were chosen} & \textbf{1st} & \textbf{2nd} & \textbf{3rd} & \textbf{Average} \\
    \hline
    At random from the full membership list & $19\%$ & $22\%$ & $19\%$ & $20\%$ \\
    From whoever arrived Tuesday evening    & $35\%$ & $36\%$ & $37\%$ & \redink{36\%} \\
    \hline
    \end{tabular}
  \end{center}
  \vspace{4pt}
  \begin{block}{Before you see the averages: which row would you trust?}
    Both rows wobble --- \textbf{variability}, unavoidable. Only the Tuesday row misses the truth
    \emph{every time}, and it misses \redink{too high}: Tuesday evening is lap-swim time, so lap
    swimmers are over-counted. \textbf{Consistent is not the same as correct} --- the tighter row
    is the wrong one.
  \end{block}
\end{frame}

% ── 1b. SAMPLING BIAS ────────────────────────────────────────────────────────
\begin{frame}[t, plain]
  \ceruleanheader{Sampling bias: who never had a chance}
  \sectionlabel{notes section 1 $\cdot$ before anyone speaks}
  \vspace{2pt}
  \begin{columns}[t]
    \begin{column}{0.47\textwidth}
      \begin{block}{Undercoverage}
        Part of the population is left \textbf{off the frame}, so those members can never be
        chosen. Tuesday evening, email-only, the social media page.
      \end{block}
    \end{column}
    \begin{column}{0.47\textwidth}
      \begin{block}{Nonresponse}
        Members are chosen properly and \textbf{never answer}. All $1{,}500$ got a card;
        $200$ mailed it back --- the frame is perfect, and the study is still broken.
      \end{block}
    \end{column}
  \end{columns}
  \vspace{12pt}
  \begin{center}
    {\large The test that separates them: \redink{could that person have been chosen?}}\\[6pt]
    {\small\color{slate} No $\Rightarrow$ undercoverage. \quad Yes, but silent $\Rightarrow$
    nonresponse. \quad Neither one shows up in your arithmetic.}
  \end{center}
\end{frame}

% ── 2a. RESPONSE BIAS ────────────────────────────────────────────────────────
\begin{frame}[t, plain]
  \ceruleanheader{Response bias: after the right person is chosen}
  \sectionlabel{notes section 2 $\cdot$ who asks, how it is worded, what order}
  \vspace{2pt}
  \begin{columns}[t]
    \begin{column}{0.5\textwidth}
      \footnotesize
      \setlength{\tabcolsep}{4pt}
      \renewcommand{\arraystretch}{1.2}
      \begin{tabular}{@{} l l @{}}
      \hline
      \textbf{Kind} & \textbf{The respondent\ldots} \\
      \hline
      Demand & answers what the study wants \\
      Social desirability & answers what looks good \\
      Dissent & says no to everything \\
      Acquiescence & agrees with everything \\
      Extreme responses & always ``strongly'' \\
      Neutral responding & always the middle box \\
      Question order & is moved by the last question \\
      \hline
      \end{tabular}
    \end{column}
    \begin{column}{0.46\textwidth}
      \begin{itemize}\setlength{\itemsep}{5pt}
        \item The lifeguard hands it out and waits; last week's complainers say ``excellent'' ---
              \redink{social desirability}.
        \item ``No opinion'' fifteen times --- \redink{neutral responding}.
        \item Agrees with two statements that contradict --- \redink{acquiescence}.
        \item Titled \emph{``Help us show the pool needs longer hours''} --- \redink{demand}.
      \end{itemize}
    \end{column}
  \end{columns}
  \vspace{6pt}
  \begin{center}
    {\small\color{slate} No sampling method prevents any of these. They happen after the right
    person has been chosen.}
  \end{center}
\end{frame}

% ── 2b. QUESTION ORDER ───────────────────────────────────────────────────────
\begin{frame}[t, plain]
  \ceruleanheader{Same words, different order}
  \sectionlabel{notes section 2 $\cdot$ question order bias}
  \vspace{0.3cm}
  \begin{center}
    \small
    \renewcommand{\arraystretch}{1.4}
    \begin{tabular}{@{} c l c @{}}
    \hline
    \textbf{Version} & \textbf{Asked first} & \textbf{Said yes to the \$40,000 lane} \\
    \hline
    A & ``Are the lap lanes overcrowded?'' & $62\%$ \\
    B & ``Should the town spend \$40,000 to add a lane?'' & $41\%$ \\
    \hline
    \end{tabular}
  \end{center}
  \vspace{0.6cm}
  \begin{center}
    {\color{cerulean}\bfseries Identical questions. Same membership. \redink{21 points apart.}}\\[8pt]
    {\small\color{slate} Ask about the crowding first and the money question inherits the
    complaint. The order of survey questions is a decision --- make it on purpose.}
  \end{center}
\end{frame}

% ── 2c. THE CRUX ─────────────────────────────────────────────────────────────
\begin{frame}[t, plain]
  \ceruleanheader{The fix that is not a fix}
  \sectionlabel[redacc]{notes section 2 $\cdot$ the whole point of today}
  \vspace{2pt}
  Ten Tuesdays instead of one: $750$ asked, $270$ say laps. \quad $270 \div 750 = {}$\redink{36\%}.
  \begin{center}
    \small
    \renewcommand{\arraystretch}{1.15}
    \begin{tabular}{@{} c c c c @{}}
    \hline
    \textbf{Members asked} & \textbf{Study says} & \textbf{Truth} & \textbf{Off by} \\
    \hline
    $75$  & $36\%$ & $20\%$ & $16$ points \\
    $750$ & $36\%$ & $20\%$ & \redink{16 points} \\
    \hline
    \end{tabular}
  \end{center}
  \textbf{Ten times the work, the same answer. Vote again.}
  \begin{block}{A larger sample makes an estimate more \emph{precise}, not more \emph{correct}}
    Size shrinks \textbf{variability}; nothing about size touches \textbf{direction}. A bigger
    biased sample is a more confident wrong answer. Only a better \emph{procedure} removes bias,
    one fix per bias: a complete frame $\to$ undercoverage; calling back the silent $\to$
    nonresponse; neutral, anonymous questions $\to$ response bias.
  \end{block}
\end{frame}

% ── WE DO — a LIVE surface: task and data only, no answers ───────────────────
\begin{frame}[t, plain]
  \ceruleanheader{Guided Practice: Cedar Ridge Apartments}
  \sectionlabel[goldacc]{We do $\cdot$ you hold the pen}
  $600$ households, \emph{every} one mailed a parking survey; $150$ cards come back, and $96$
  say parking is a problem. Management is about to ask the owners for a bigger lot.
  \begin{itemize}\setlength{\itemsep}{4pt}
    \item[(a)] What share answered? What share of \emph{those} complained?
    \item[(b)] Every household got a card. Which bias --- and could a silent household have been chosen?
    \item[(c)] Scale $64\%$ to $600$; the census says $240$ of $600$. Off by how much, which way?
    \item[(d)] Twice the cards, $64\%$ again; the car owners answered at twice the rate. Why didn't
          more cards help?
    \item[(e)] \emph{``Don't you agree the lot is dangerously overcrowded?''} Name the bias; rewrite it neutral.
  \end{itemize}
  \begin{block}{The question behind the question}
    If your answer to (d) is ``they still need more cards,'' you are still at the hook vote.
  \end{block}
\end{frame}

% ── DEBRIEF ──────────────────────────────────────────────────────────────────
\begin{frame}[t, plain]
  \ceruleanheader{Debrief: more of a biased sample is still biased}
  \sectionlabel{10 minutes $\cdot$ pens down, papers up --- we say these aloud}
  \vspace{4pt}
  \begin{enumerate}\setlength{\itemsep}{5pt}
    \item The ten Tuesdays back on the board. Say the sentence: \emph{$750$ is not closer than
          $75$, because size shrinks variability and this study's problem is direction.}
    \item The $200$ cards that came back. Could a silent card holder have been chosen? Then
          which bias is it?
    \item Cedar Ridge (d): the direction and the reason --- \emph{too high, because}\,\ldots
  \end{enumerate}
  \vspace{6pt}
  \begin{block}{Cold check --- hands down, thirty seconds}
    Bayside Middle School, $750$ students. The principal asks $30$ students in the library at
    lunch whether lunch is too short: $24$ say yes ($80\%$). Told the sample was too small, she
    asks $60$ library students: $48$ say yes --- $80\%$ again. \textbf{Name the bias in her plan.
    Did the larger sample fix it? Why not?}
  \end{block}
\end{frame}

% ── YOU DO — the homework is the individual practice, started in class ───────
\begin{frame}[t, plain]
  \ceruleanheader{You do: start the homework}
  \sectionlabel[goldacc]{Last 10 minutes $\cdot$ on your own $\cdot$ finish it at home}
  \begin{itemize}\setlength{\itemsep}{5pt}
    \item \textbf{Millbrook Public Library, two studies and two answers} --- a front-desk pile
          that says $80\%$ and a proper phone sample that says $45\%$.
    \item Population, sample, who chose it; both percents and both predictions; the gap in
          points and in card holders.
    \item \textbf{Item 4:} four weeks at the desk versus the phone sample --- which is bigger,
          which is closer, and what the $360$ would have said.
    \item Four response biases to name; a $21$-point order effect; a random draw where only $78$
          of $120$ pick up.
  \end{itemize}
  \begin{block}{Silent and alone --- I am circulating. Packet pages, scored out of 12, due at the start of the first class after two study halls.}
    Start with \textbf{1, 4, and 6}. Item 4 is the one I am reading over your shoulder --- if
    your answer to ``what would the $360$ say'' is anywhere near $45\%$, flag me before you
    leave.
  \end{block}
\end{frame}

% ── CLOSE ────────────────────────────────────────────────────────────────────
{
\setbeamercolor{background canvas}{bg=cerulean}
\begin{frame}[plain]
  \vspace{1.5cm}\hspace{0.8cm}%
  \begin{minipage}{0.86\textwidth}
    {\color{goldacc}\bfseries\small BEFORE YOU GO}\\[0.7cm]
    {\color{white}\bfseries\Large Bias is a direction, not bad luck. A larger sample makes an
    estimate more precise, not more correct --- only a better procedure removes the lean. Ask
    who could never have been chosen, who stayed silent, and what the question told them to
    say.}\\[0.9cm]
    {\color{frostmid}\small Next: Lesson 1.5 --- every study in this unit so far has only
    \emph{watched}. Observational studies: what they can and cannot prove.}
  \end{minipage}
\end{frame}
}

\end{document}
